% generated by b1_ailoop.py; do not edit
\newcommand{\BoneControl}{100}\newcommand{\BoneRamp}{200}
\newcommand{\BoneFA}{0}
\newcommand{\BoneEarly}{193}\newcommand{\BoneEarlyShare}{96}\newcommand{\BoneLead}{13.0}
\newcommand{\BoneGain}{17.0}\newcommand{\BonePos}{100}\newcommand{\BoneNeg}{0}
\newcommand{\BoneBelow}{200}\newcommand{\BoneBelowN}{200}
\newcommand{\BoneShareCrisis}{0.46}\newcommand{\BoneStatic}{0.82}
\newcommand{\BonePone}{passed}\newcommand{\BonePtwo}{passed}\newcommand{\BonePthree}{passed}\newcommand{\BonePfour}{passed}

\paragraph{Design.} The market is the closed loop of the thesis on five-minute bars, with 50 agents. Independent learners each have their own memory and learn their response online; with reverting prices they learn to trade against moves. A vendor model is a pre-trained momentum product with one memory rate; every agent that adopts it responds identically. Adoption ramps from none to 80\% of the agents between days 24 and 96 of a 120-day run, in two versions: a fast vendor model (memory rate $0.9$ per bar) and a slow one ($0.995$, a half-life of about eleven hours); in the control nobody adopts. The observer sees only the price displacement and the anonymous flow and runs the pooled estimator of Theorem~\ref{thm:pooled} daily on a ten-day window, alone (standard) or with the flow equation augmented by the signal at a disclosed memory bound of $0.995$ (disclosed). A crisis is the first day on which the displacement exceeds ten times its standard deviation before adoption. The warning rules were fixed in an exploratory pilot on other seeds: the certificate of a positive margin is lost; the certificate of a margin above a buffer of $0.1$ is lost; or the projected time to the boundary, the estimate divided by its rate of decline over five days, falls below ten days. Test~B1 was registered, with its code and four primaries, before its seeds (1001--1100, three scenarios, 300 runs) were run, and it ran once.

\paragraph{Result.} All four primaries passed. In the \BoneControl{} control runs the standard estimator raised no false alarm under any rule (\BoneFA{} for the time-to-boundary rule). Every one of the \BoneRamp{} runs with adoption ended in a crisis, and the time-to-boundary rule warned before it in \BoneEarly{} ($\BoneEarlyShare\%$), with a median lead of \BoneLead{} days (\BoneLeadFast{} for the fast vendor model, \BoneLeadSlow{} for the slow one). The level of the margin alone gave no warning: the standard certificate was lost on the day of the crisis (median lead \BoneLostLead{} days). With the slow vendor model, the disclosed estimator's buffer warning came earlier than the standard one in \BonePos{} of 100 runs and later in \BoneNeg{}, a median gain of \BoneGain{} days (median leads \BoneBufDiscSlow{} and \BoneBufStdSlow{} days). And in all \BoneBelowN{} runs the crisis came at an adoption share below the static threshold of Proposition~\ref{prop:monoculture}, computed with the learners frozen when adoption began: a median share of \BoneShareCrisis{} against a median threshold of \BoneStatic{}.

\paragraph{What the secondaries show.} Disclosure has a price in false alarms. The disclosed estimator is noisier: in the controls its buffer rule fired in \BoneFAdiscBuf{} of 100 runs, its lost-certificate rule in \BoneFAdiscLost{}, and its time-to-boundary rule in \BoneFAdiscTtz{}, so only the standard estimator's trend and the disclosed estimator's buffer are usable together, and the gain of Primary~3 is earned partly by a rule that would false-alarm about one time in six. The standard estimator's daily estimate differed from the true margin by a median of \BoneTrack{}, mostly as a lag of a few days behind a margin that falls quickly.

\paragraph{Reading.} Three findings matter for supervision. The certificate's level is not an early warning, but its trend is: a margin falling toward zero is visible in public data about two weeks before the crisis, with no false alarms in markets where nothing changes. A disclosed memory bound buys more warning when the shared model is slow, which is the case Theorem~\ref{thm:hidden} says prices alone cannot see, at a cost in false alarms that a supervisor would have to weigh. And learning makes monoculture dangerous earlier than static analysis predicts: as the adopters push prices, the independent learners stop trading against the move and start following it, so the market crosses the boundary at about half the adoption share at which a vendor-by-vendor assessment, freezing everyone else, would put it. That last finding is the strongest argument for supervising the loop rather than the models: the danger is created by the interaction, and an assessment of any one model, however careful, misses it.
